Predicting drug--target interactions (DTIs) from molecular and protein information can prioritize candidate mechanisms, but apparent accuracy can depend on what information crosses an evaluation split. We developed ArnoldiGCL-S1/S2, a graph-based DTI predictor that combines separate drug and protein projections with seven similarity-neighbour descriptors computed only from positive training-fold interactions. Across BioSNAP, Human and BindingDB under random, cold-drug, cold-protein and cold-pair splits, the full model had the highest mean AUPR among the principal baselines in 7/12 prespecified conditions and exceeded its structural-only variant in 11/12 conditions. The latter gains ranged from 0.013 to 0.396 AUPR; 11 descriptive paired contrasts had Benjamini--Hochberg-adjusted $p<0.05$. The model was not the highest-AUPR principal baseline in BioSNAP Cold protein, Human Cold protein, Human Cold pair, BindingDB Cold protein, BindingDB Cold pair, demonstrating that the side information does not confer a uniform advantage. These results support leakage-safe similarity-neighbour summaries as a useful, but split-dependent, source of evidence for computational DTI prioritization; they do not establish prospective or experimental target binding.
